% GENERATED FILE. Edit canonical lessons, then run npm run generate:books.
% canonical-source-hash: fnv1a64:1d3c08eb21677fbb
% canonical-lessons: ES-C440-boda, ES-C440-cumpleanos, ES-C440-celebrar, ES-C440-disfrutar, ES-C440-humor, ES-R440-repaso-fiesta, ES-C440-sintesis-fiesta

\chapter{La Fiesta Y El Año}
\label{ch:es-fiesta}
\hlchaptermodality{440}

\begin{chapteropening}
\textbf{By the end of this chapter:} \emph{I can be invited to a wedding or a birthday, say I went, say what mood I was in, and tell a notice's \textquotedblleft{}is held\textquotedblright{} apart from a friend's \textquotedblleft{}we're celebrating\textquotedblright{}.}

Read the message or notice this chapter's words exist for, then say the same thing back.
\end{chapteropening}

\section[la boda]{la boda — the vows, taken as one thing}

\label{lesson:ES-C440-boda}



\begin{quote}
Say \textbf{la fiesta}, \textbf{la iglesia} and \textbf{la familia}. One day puts all three in the same place. Name it.
\end{quote}

\subsection*{You'll want to know: la boda}
\textbf{la boda} — the wedding.

\emph{Vamos a una boda el sábado} — we are going to a wedding on Saturday. Spanish uses the singular for the event and the plural \textbf{las bodas} for an anniversary of it: \emph{las bodas de plata} are a silver wedding, twenty-five years on.

The \textbf{b} is soft between vowels — \emph{la BO-da}, with the lips barely closing.

\begin{cousinweb}
\textbf{boda} looks like one solid piece and is not. It is Latin \textbf{vōta}, \emph{the vows} — the plural of \textbf{vōtum}, a vow.

A wedding was named for what is said at it, and the name is plural because more than one vow is spoken. Then Spanish stopped hearing the plural and treated the whole word as a single feminine noun, which is why you can now say \emph{una boda}, one-of-the-vows, without it sounding odd to anybody.

English kept the same root in \textbf{vote}, \textbf{votive} and \textbf{devotion}, all of them about something pledged.
\end{cousinweb}

\begin{culture}
This is not a one-off. Latin had a third gender, \textbf{neuter}, whose plural ended in \textbf{-a}, and Spanish lost that gender entirely. The plurals were left standing with a feminine-looking ending and nothing to belong to, so the language took them as feminine singulars.

\noindent
\begin{tabularx}{\linewidth}{@{}>{\raggedright\arraybackslash}X>{\raggedright\arraybackslash}X@{}}
\toprule
\textbf{Latin, a neuter plural} & \textbf{Spanish, a feminine singular} \\
\midrule
vōta, the vows & \textbf{la boda} \\
folia, the leaves & la hoja \\
ligna, the firewoods & la leña \\
\bottomrule
\end{tabularx}

Every one of those is a \textbf{countable plural read as a single mass or event}, and that is the shape to look for: a Spanish feminine noun naming one thing made of many. It is a third way for a word to hide its pieces — not a false ending, not a false prefix, but a number the language forgot.

\texttt{cerrado} in the last chapter came apart the moment you looked. \textbf{boda} will not, and no amount of looking at modern Spanish would tell you why.
\end{culture}

\subsection*{Your turn}
Say these aloud:

\begin{itemize}
\raggedright
  \item \emph{la boda}
  \item we are going to a wedding on Saturday
  \item what \emph{boda} originally meant, and why it was plural
\end{itemize}

\begin{tcolorbox}[breakable,colback=teal!4,colframe=teal!35!black,title={Before you move on}]
The wedding? (\textbf{La boda}.) What was the Latin word? (\emph{\textbf{Vōta}}, the vows.) What number was it? (\textbf{Plural}.)
\end{tcolorbox}
\section[el cumpleaños]{el cumpleaños — the years you have completed}

\label{lesson:ES-C440-cumpleanos}



\begin{quote}
Say \emph{¿Cuántos años tienes?} Say \textbf{el aniversario}. Say \textbf{la boda}. One day a year answers the first question, and it is not called an \emph{aniversario}.
\end{quote}

\subsection*{You'll want to know: el cumpleaños}
\textbf{el cumpleaños} — the birthday.

It ends in \textbf{-s} and it is \textbf{singular}: \emph{el cumpleaños}, \emph{los cumpleaños}. The word does not change; only the article tells you how many.

\emph{¡Feliz cumpleaños!} — happy birthday. \emph{Mañana es mi cumpleaños} — tomorrow is my birthday.

\begin{cousinweb}
Two words fused into one: \textbf{cumple} plus \textbf{años}.

\textbf{cumplir} is to complete, to fulfil, to carry something through to its end. \emph{Cumplo treinta} — I complete thirty. So a \textbf{cumple-años} is the day the years get completed, and the \textbf{-s} you can hear is the plural of \emph{años} trapped inside a word that is now singular.

That is why \emph{el cumpleaños} looks plural and is not. It is not an exception to be memorised; it is a sentence that hardened.
\end{cousinweb}

\begin{culture}
Spanish has two words for the same square on the calendar, and they disagree about what is happening.

\noindent
\begin{tabularx}{\linewidth}{@{}>{\raggedright\arraybackslash}X>{\raggedright\arraybackslash}X@{}}
\toprule
\textbf{the word} & \textbf{what it says is happening} \\
\midrule
el aniversario & the \textbf{year turns} \\
el cumpleaños & \textbf{you complete} years \\
\bottomrule
\end{tabularx}

\emph{Aniversario} is \emph{annus} plus \emph{vertere}, the year coming round again — the calendar does the work, and you stand still. \emph{Cumpleaños} puts the work on you: the years are something you finish, one at a time, like a task.

Spanish keeps both and uses them for different things. A wedding gets an \textbf{aniversario} because a date recurs; a person gets a \textbf{cumpleaños} because a person is doing the ageing. That is not a rule anyone wrote down — it is the metaphor inside each word still deciding where it is welcome.
\end{culture}

\subsection*{Your turn}
Say these aloud:

\begin{itemize}
\raggedright
  \item \emph{el cumpleaños}, then \emph{los cumpleaños}
  \item tomorrow is my birthday
  \item why a wedding has an \emph{aniversario} and a person has a \emph{cumpleaños}
\end{itemize}

\begin{tcolorbox}[breakable,colback=teal!4,colframe=teal!35!black,title={Before you move on}]
The birthday? (\textbf{El cumpleaños}.) Is it plural? (\textbf{No — singular, with a plural trapped inside}.) Which two words is it made of? (\emph{\textbf{Cumple}} and \emph{\textbf{años}}.)
\end{tcolorbox}
\section[celebrar]{celebrar — to make a crowd of it}

\label{lesson:ES-C440-celebrar}



\begin{quote}
Say \textbf{el cumpleaños}, \textbf{la fiesta} and \textbf{la música}. You have the day, the party and the sound. Name the verb that puts them together.
\end{quote}

\subsection*{You'll want to know: celebrar}
\textbf{celebrar} — to celebrate.

\emph{Celebramos su cumpleaños en casa} — we are celebrating her birthday at home. \emph{Celebrarse} is the form a notice uses: \emph{la reunión se celebra el martes} — the meeting is held on Tuesday.

That second use is worth having early. On an official notice, \textbf{celebrarse} does not mean anybody is enjoying themselves; it means the event \textbf{takes place}.

\begin{cousinweb}
Latin \textbf{celebrāre}, from \textbf{celeber} — \emph{crowded, much-visited, thronged}.

A \emph{celeber locus} was a place people went to in numbers. To \textbf{celebrate} something was therefore to \textbf{make a crowd of it}, and only later to enjoy it. The idea underneath is attendance, not happiness.

English took the same root twice: \textbf{celebrate}, and \textbf{celebrity} — a person who has a crowd.
\end{cousinweb}

\begin{culture}
That original meaning is exactly why a town-hall notice can say \emph{se celebra} for an event nobody would call festive.

\noindent
\begin{tabularx}{\linewidth}{@{}>{\raggedright\arraybackslash}X>{\raggedright\arraybackslash}X@{}}
\toprule
\textbf{what the sign says} & \textbf{what it means} \\
\midrule
se celebra la fiesta & the party takes place \\
se celebra la reunión & the meeting takes place \\
\bottomrule
\end{tabularx}

Both are \emph{celebrar}, because both are a gathering. English would need two verbs here, and a reader who knows only the happy sense will misread the second one as a joke.

When you meet \textbf{celebrarse} on a notice, read it as \textbf{is held}. The crowd is what the word is really about, and a meeting has one.
\end{culture}

\subsection*{Your turn}
Say these aloud:

\begin{itemize}
\raggedright
  \item \emph{celebrar}
  \item we are celebrating my birthday at home
  \item what \emph{la reunión se celebra el martes} is telling you
\end{itemize}

\begin{tcolorbox}[breakable,colback=teal!4,colframe=teal!35!black,title={Before you move on}]
To celebrate? (\textbf{Celebrar}.) What did \emph{celeber} mean? (\textbf{Crowded}.) What does \emph{se celebra} mean on a notice? (\textbf{It is held}.)
\end{tcolorbox}
\section[disfrutar]{disfrutar — to take the fruit of it}

\label{lesson:ES-C440-disfrutar}



\begin{quote}
Say \textbf{la fruta}, \textbf{descansar} and \textbf{celebrar}. Two of those three share a prefix shape. Hold on to that.
\end{quote}

\subsection*{You'll want to know: disfrutar}
\textbf{disfrutar} — to enjoy.

\emph{Disfruté mucho de la boda} — I really enjoyed the wedding. \emph{¡Que disfrutes!} — enjoy yourself, said as you hand somebody a ticket or wave them off.

It usually takes \textbf{de}: \emph{disfrutar \textbf{de} algo}. Dropping the \emph{de} is common in speech and still reads as informal in writing.

\begin{cousinweb}
\textbf{dis-} plus \textbf{frut-}, and the second half is \textbf{la fruta} you already have. Both come from Latin \textbf{frūctus}, which meant \emph{produce, yield, the thing a plant gives you}.

To \textbf{disfrutar} something is to \textbf{take the fruit of} it — to get the yield. English has the same picture in \emph{to enjoy the fruits of your labour}, which most speakers hear as a metaphor and which Spanish has compressed into one verb.
\end{cousinweb}

\begin{culture}
Now the trap, and it corrects the test you learned two chapters ago.

\textbf{descansar} is \textbf{des-} plus \emph{cansar}, and the prefix \textbf{undoes}: it takes the tiredness away. It would be natural to read \textbf{dis-} in \emph{disfrutar} the same way and land on \emph{to un-fruit}, which would be the opposite of what the word means.

\noindent
\begin{tabularx}{\linewidth}{@{}>{\raggedright\arraybackslash}X>{\raggedright\arraybackslash}X@{}}
\toprule
\textbf{word} & \textbf{what the prefix does} \\
\midrule
descansar & reverses — removes the tiredness \\
\textbf{disfrutar} & \textbf{intensifies — takes the yield fully} \\
\bottomrule
\end{tabularx}

The piece-test still works: take \textbf{dis-} off and \emph{frutar} stands, so the join is real. What the test does \textbf{not} tell you is what the join \emph{means}. A prefix that usually reverses does not always reverse, and the only way to know is the word's sense, not its shape.

That is the honest limit of taking words apart. It shows you the seams; it does not read them for you.
\end{culture}

\subsection*{Your turn}
Say these aloud:

\begin{itemize}
\raggedright
  \item \emph{disfrutar}
  \item I really enjoyed the wedding we celebrated
  \item why \emph{dis-} in \emph{disfrutar} is not the \emph{des-} of \emph{descansar}
\end{itemize}

\begin{tcolorbox}[breakable,colback=teal!4,colframe=teal!35!black,title={Before you move on}]
To enjoy? (\textbf{Disfrutar}.) Which preposition follows it? (\emph{\textbf{De}}.) What does the second half mean? (\textbf{Fruit, yield}.)
\end{tcolorbox}
\section[el humor]{el humor — the fluid you are in}

\label{lesson:ES-C440-humor}



\begin{quote}
Say \textbf{reír}, \textbf{alegre} and \textbf{disfrutar}. All three describe a person having a good time. Now name the thing that person is \textbf{in}.
\end{quote}

\subsection*{You'll want to know: el humor}
\textbf{el humor} — mood, and also humour.

\emph{Estoy de buen humor} — I am in a good mood. \emph{Está de mal humor} — he is in a bad mood. \emph{Tiene mucho sentido del humor} — she has a great sense of humour.

The \textbf{h} is silent here as everywhere in Spanish: \emph{el u-MOR}. And note the preposition — Spanish says you are \textbf{de} a mood, of it, not \emph{in} it.

\begin{cousinweb}
Latin \textbf{ūmor} meant one thing only: \textbf{moisture, a fluid}. English keeps that original sense in \textbf{humid}.

Medieval medicine held that a body ran on four fluids, and that whichever one you had most of set your temperament. Your \emph{humour} was literally which liquid you were full of. The theory is long dead, and the word survived it — so when a Spanish speaker says \emph{estar de buen humor}, the sentence still says \emph{to be of a good fluid}, and nobody notices.
\end{cousinweb}

\begin{culture}
One Spanish word covers what English splits into \textbf{mood} and \textbf{humour}, and the split happened in English rather than the join happening in Spanish.

\noindent
\begin{tabularx}{\linewidth}{@{}>{\raggedright\arraybackslash}X>{\raggedright\arraybackslash}X@{}}
\toprule
\textbf{Spanish} & \textbf{English needs} \\
\midrule
estar de buen humor & to be in a good \textbf{mood} \\
sentido del humor & sense of \textbf{humour} \\
\bottomrule
\end{tabularx}

English narrowed \emph{humour} onto the funny sense and borrowed \emph{mood} from its own older stock to carry the rest. Spanish never divided the work, so \emph{el humor} still holds both — the state you are in, and the thing that makes other people laugh.

That is worth noticing whenever an English pair maps to one Spanish word. The usual cause is not that Spanish is vaguer; it is that English acquired a second word from somewhere and gave it half the job.
\end{culture}

\subsection*{Your turn}
Say these aloud:

\begin{itemize}
\raggedright
  \item \emph{el humor}
  \item I am in a good mood and I enjoyed it a lot
  \item \emph{sentido del humor}, and what the word first meant
\end{itemize}

\begin{tcolorbox}[breakable,colback=teal!4,colframe=teal!35!black,title={Before you move on}]
Mood? (\textbf{El humor}.) Which preposition goes with it? (\emph{\textbf{De}}.) What did the Latin word mean? (\textbf{Moisture}.)
\end{tcolorbox}
\section[(boda, cumpleaños, celebrar …]{(boda, cumpleaños, celebrar, disfrutar, humor) — from cold}

\label{lesson:ES-R440-repaso-fiesta}



\begin{quote}
Close the lessons before this one. Say all five: \emph{the wedding}, \emph{the birthday}, \emph{to celebrate}, \emph{to enjoy}, \emph{the mood}.
\end{quote}

\begin{grammarlens}[title={Grammar Lens: three ways a word hides its pieces}]
You now have a small collection of them, and they fail differently.

\noindent
\begin{tabularx}{\linewidth}{@{}>{\raggedright\arraybackslash}X>{\raggedright\arraybackslash}X@{}}
\toprule
\textbf{the word} & \textbf{what is hidden} \\
\midrule
\textbf{boda} & a \textbf{plural} the language forgot \\
\textbf{cumpleaños} & a whole \textbf{phrase}, fused \\
\textbf{disfrutar} & a real join whose \textbf{meaning} misleads \\
\bottomrule
\end{tabularx}

\emph{Boda} cannot be recovered from modern Spanish at all — only Latin says why it is one word. \emph{Cumpleaños} can: say it slowly and \emph{cumple} and \emph{años} fall apart in your mouth. \emph{Disfrutar} comes apart cleanly and then lies about what the prefix is doing.

The piece-test tells you \textbf{whether} there is a join. It never tells you what the join means, and \emph{disfrutar} is the word that proves it.
\end{grammarlens}

\subsection*{Your turn}
Say these aloud:

\begin{itemize}
\raggedright
  \item we are celebrating my birthday on Saturday
  \item I am in a good mood
  \item what \emph{se celebra} means on an official notice
\end{itemize}

\begin{grammarlens}[title={What you've built}]
\textbf{You can now be invited to something, say you went, and say how it was.} That is the whole social round-trip, and it is what a message from a friend actually contains.

You also have the pair that will keep paying: \textbf{celebrarse} on a notice means \emph{is held}, not \emph{is enjoyed}, and \textbf{el humor} is one word where English keeps two.
\end{grammarlens}

\begin{tcolorbox}[breakable,colback=teal!4,colframe=teal!35!black,title={Before you move on}]
The birthday? (\textbf{El cumpleaños}.) The wedding? (\textbf{La boda}.) In a good mood? (\textbf{De buen humor}.) Which of those two nouns hides a plural? (\emph{\textbf{Boda}} — \emph{vōta}, the vows.)
\end{tcolorbox}
\section[Practice]{(una invitación y una respuesta) — read it, then do it yourself}

\label{lesson:ES-C440-sintesis-fiesta}



\begin{quote}
Say \emph{la boda} and \emph{celebrar}. The same verb is about to appear twice in front of you, meaning two different things.
\end{quote}

\subsection*{Your turn}
Read both, then answer.

\textbf{1 — un mensaje}

\begin{quote}
¡Hola! El sábado \textbf{celebramos} mi \textbf{cumpleaños} en casa de mis padres. Nada especial, solo familia y unos amigos. ¿Te apuntas?
\end{quote}

\textbf{2 — un aviso en el portal}

\begin{quote}
La junta de vecinos se \textbf{celebrará} el jueves 9 a las 19:00 en el portal. Se ruega puntualidad.
\end{quote}

Say these aloud:

\begin{itemize}
\raggedright
  \item which of the two is something you would enjoy
  \item what \emph{se celebrará} is telling you in the second one
  \item where and when each one happens
\end{itemize}

\begin{culture}
The two uses look identical and are told apart by one thing: whether the verb is reflexive.

\textbf{Celebramos mi cumpleaños} — \emph{we} celebrate, somebody is doing the enjoying. \textbf{Se celebrará la junta} — it \emph{gets} held, and no one is named, because a notice does not care who convenes a meeting. The reflexive drains the person out of the sentence, and what is left is the bare fact that the thing takes place.

So the answer to \emph{which of the two would you enjoy} is not a matter of tone. It is in the grammar: the message has a \textbf{we}, and the notice does not.

Now answer the message the way a friend would: say yes, say what mood you are in, and say you will enjoy it.
\end{culture}

\begin{tcolorbox}[breakable,colback=teal!4,colframe=teal!35!black,title={Before you move on}]
Now your turn, out loud and then in writing. Four sentences: accept the invitation, say what mood you are in, say you are sure you will enjoy it, and mention one wedding or birthday of your own coming up.

Then check yourself: did you use \textbf{de} after \emph{humor} and after \emph{disfrutar}?
\end{tcolorbox}
